% Source: time_series_multiplicative_seasonality.tex
\begin{notechapter}
% ---- chapter-local definitions ----
\lstdefinestyle{mystyle}{
    backgroundcolor=\color{white},
    commentstyle=\color{green},
    keywordstyle=\color{blue},
    numberstyle=\tiny\color{gray},
    stringstyle=\color{purple},
    basicstyle=\ttfamily\footnotesize,
    breakatwhitespace=false,
    breaklines=true,
    captionpos=b,
    keepspaces=true,
    numbers=left,
    numbersep=5pt,
    showspaces=false,
    showstringspaces=false,
    showtabs=false,
    tabsize=4
}
\lstset{style=mystyle}
% -----------------------------------
\chapter{Time Series with Multiplicative Seasonality}\label{ch:multiplicative-seasonality}

\section{Time Series with Multiplicative Seasonality}

A time series with multiplicative seasonality can be modeled as:
\[
Y_t = T_t \times S_t \times \epsilon_t
\]
where:
\begin{itemize}
    \item \( Y_t \) is the observed value at time \( t \),
    \item \( T_t \) is the trend component,
    \item \( S_t \) is the seasonal component,
    \item \( \epsilon_t \) is the random error (assumed negligible).
\end{itemize}

For additive seasonality, the model is:
\[
Y_t = T_t + S_t + \epsilon_t
\]

\section{Constructing the Time Series}

\subsection{Trend Component (\( T_t \))}
The trend component represents the long-term growth or decline in the data.
In this example, we assume a linear trend:
\[
T_t = 100 + 4t
\]
where \( t \) is the time period (in months).
This means the sales increase by 4 units every month.

\subsection{Seasonal Component (\( S_t \))}
The seasonal component captures periodic fluctuations.
For multiplicative seasonality, we use a sinusoidal function scaled
to create a multiplicative effect:
\[
S_t = 1 + 0.5 \sin\left(\frac{2\pi t}{12}\right)
\]
This creates a seasonal pattern that repeats every 12 months (1 year).
Because $S_t$ multiplies the trend, the seasonal swings ($\pm 50\%$ of the trend here, since the amplitude is $0.5$) grow proportionally with the level of the series. Note also that $S_t$ averages to $1$ over a full period.

\subsection{Observed Series (\( Y_t \))}
The observed series is obtained by multiplying the trend and seasonal
components:
\[
Y_t = T_t \times S_t
\]
Substituting the expressions for \( T_t \) and \( S_t \):
\[
Y_t = (100 + 4t) \times \left(1 + 0.5 \sin\left(\frac{2\pi t}{12}\right)\right)
\]

\subsection{Example Data}
The table below shows the observed data \( Y_t \) for 36 months (3 years), with $t=1$ in January of year 1 (values rounded to one decimal). Note how the seasonal swings grow with the level of the series: this is the signature of multiplicative seasonality.

\begin{center}
\begin{tabular}{|c|c|c|c|}
\hline
Month & Year 1 Sales & Year 2 Sales & Year 3 Sales \\
\hline
January & 130.0 & 190.0 & 250.0 \\
February & 154.8 & 223.5 & 292.3 \\
March & 168.0 & 240.0 & 312.0 \\
April & 166.2 & 235.0 & 303.8 \\
May & 150.0 & 210.0 & 270.0 \\
June & 124.0 & 172.0 & 220.0 \\
July & 96.0 & 132.0 & 168.0 \\
August & 74.8 & 102.1 & 129.3 \\
September & 68.0 & 92.0 & 116.0 \\
October & 79.4 & 106.6 & 133.8 \\
November & 108.0 & 144.0 & 180.0 \\
December & 148.0 & 196.0 & 244.0 \\
\hline
\end{tabular}
\end{center}

\section{Isolating the Trend Using Moving Averages}

To isolate the trend from the observed series \( Y_t \), we can use a
moving average. A moving average smooths out short-term fluctuations
and highlights the long-term trend. For monthly data with a 12-month
seasonality, a 12-month moving average is appropriate.

The moving average \( \text{MA}_t \) at time \( t \) is calculated as:
\[
\text{MA}_t = \frac{1}{12} \sum_{k=t-6}^{t+5} Y_k
\]
which is the convention used in the code below. Since the window length is even, this average is centered at $t - \tfrac{1}{2}$ rather than at $t$. Classical decomposition therefore uses the \emph{centered} $2\times 12$ moving average
\[
\text{MA}_t = \frac{1}{24}\Big(Y_{t-6} + 2\sum_{k=t-5}^{t+5} Y_k + Y_{t+6}\Big),
\]
which is exactly centered at $t$ and still removes a seasonal pattern of period 12.

\subsection{Convolution for Moving Average}
The moving average is computed efficiently using \textbf{convolution}.
Convolution is a mathematical operation that combines two sequences to
produce a third sequence. In this case, we convolve the time series
\( Y_t \) with a kernel of weights.

For a 12-month moving average, the kernel is:
\[
\text{kernel} = \left[\frac{1}{12}, \frac{1}{12}, \dots, \frac{1}{12}\right]
\]
This kernel is applied to the time series \( Y_t \) using the
`np.convolve` function. The convolution operation is performed as
follows:
\[
\text{moving\_avg} = Y_t \ast \text{kernel}
\]
where \( \ast \) denotes the convolution operation.

The `mode='valid'` argument ensures that the convolution is computed
only where the kernel fully overlaps with the time series.
This avoids edge effects and results in a smoothed series of
length \( N - 11 \), where \( N \) is the length of \( Y_t \).

\section{Removing the Trend}

Once the trend \( \text{MA}_t \) is isolated, we can remove it from
the observed series \( Y_t \) to analyze the seasonal and residual
components. For multiplicative seasonality, this is done by dividing
the observed series by the trend:
\[
\text{Detrended Series} = \frac{Y_t}{\text{MA}_t}
\]

\section{Python Code for Trend Isolation and Visualization}

The following Python code generates the time series, isolates the trend
using a moving average, removes the trend, and plots the results.

\begin{pyblock}
import numpy as np
import matplotlib.pyplot as plt

# Time periods (36 months)
t = np.arange(1, 37)

# Trend component (T_t)
T_t = 100 + 4 * t

# Seasonal component (S_t)
S_t = 1 + 0.5 * np.sin(2 * np.pi * t / 12)

# Observed series (Y_t)
Y_t = T_t * S_t

# Calculate the 12-month moving average to isolate the trend
window_size = 12
moving_avg = np.convolve(Y_t, np.ones(window_size)/window_size, mode='valid')

# Adjust the time axis for the moving average
t_moving_avg = t[window_size//2 : len(t) - window_size//2 + 1]

# Remove the trend by dividing Y_t by the moving average
detrended_series = Y_t[window_size//2 : len(t) - window_size//2 + 1] / moving_avg

# Plot 1: Time series and moving average
plt.figure(figsize=(10, 6))
plt.plot(t, Y_t, marker='o', label="Observed Data")
plt.plot(t_moving_avg, moving_avg, linestyle='--', label="12-Month Moving Average (Trend)")
plt.title("Time Series with Multiplicative Seasonality and Trend Isolation")
plt.xlabel("Months")
plt.ylabel("Sales")
plt.xticks(np.arange(1, 37, 3), 
           ["Jan Y1", "Apr Y1", "Jul Y1", "Oct Y1", 
            "Jan Y2", "Apr Y2", "Jul Y2", "Oct Y2", 
            "Jan Y3", "Apr Y3", "Jul Y3", "Oct Y3"])
plt.grid(True)
plt.legend()
plt.savefig("time_series_with_trend.png")  # Save the plot
plt.close()  # Close the figure

# Plot 2: Detrended series
plt.figure(figsize=(10, 6))
plt.plot(t_moving_avg, detrended_series, marker='o', label="Detrended Data")
plt.title("Detrended Time Series (Seasonal and Residual Components)")
plt.xlabel("Months")
plt.ylabel("Sales (Detrended)")
plt.xticks(np.arange(1, 37, 3), 
           ["Jan Y1", "Apr Y1", "Jul Y1", "Oct Y1", 
            "Jan Y2", "Apr Y2", "Jul Y2", "Oct Y2", 
            "Jan Y3", "Apr Y3", "Jul Y3", "Oct Y3"])
plt.grid(True)
plt.legend()
plt.savefig("detrended_series.png")  # Save the plot
plt.close()  # Close the figure
\end{pyblock}

\section{Generated Plots}

\subsection{Plot 1: Time Series with Trend Isolation}
The first plot shows the observed data \( Y_t \) and the 12-month
moving average (trend).

\begin{figure}[!h]
    \centering
    \includegraphics[width=0.8\textwidth]{time_series_with_trend.png}
    \caption{Time Series with Multiplicative Seasonality and Trend Isolation}
\end{figure}

\subsection{Plot 2: Detrended Time Series}
The second plot shows the detrended series, which represents the
seasonal and residual components after removing the trend.

\begin{figure}[!h]
    \centering
    \includegraphics[width=0.8\textwidth]{detrended_series.png}
    \caption{Detrended Time Series (Seasonal and Residual Components)}
\end{figure}

\section{Explanation of the Detrended Series}

The detrended series is obtained by dividing the observed series
\( Y_t \) by the moving average \( \text{MA}_t \). This removes the
trend component, leaving behind the seasonal and residual components:
\[
\text{Detrended Series} = \frac{Y_t}{\text{MA}_t}
\]
The detrended series highlights the seasonal patterns and any remaining
random fluctuations.
\end{notechapter}
